% BHU PYQ -- Manually transcribed & error-corrected
% Paper: BCF-299: Financial Analysis
% Exam: B.Com. (Hons.) F.M.M. Semester IV Examination, 2015-16

\documentclass[11pt,a4paper]{article}
\usepackage[a4paper,top=2.5cm,bottom=2.5cm,left=2.8cm,right=2.8cm,headheight=14pt]{geometry}
\usepackage{amsmath,amssymb}
\usepackage[T1]{fontenc}
\usepackage[utf8]{inputenc}
\usepackage{lmodern,microtype,fancyhdr,enumitem,hyperref,lastpage,parskip}

\hypersetup{
    colorlinks=true,
    urlcolor=black,
    linkcolor=black,
    pdftitle={BCF-299: Financial Analysis -- BHU B.Com. (Hons.) FMM Sem IV 2015-16}
}

\pagestyle{fancy}
\fancyhf{}
\renewcommand{\headrulewidth}{0.4pt}
\renewcommand{\footrulewidth}{0.4pt}
\fancyhead[L]{\small   \textit{Banaras Hindu University}}
\fancyhead[R]{\small   \textit{BCF-299: Financial Analysis}}
\fancyfoot[C]{\small Page  \thepage\ of \pageref{LastPage}}


\newlist{parts}{enumerate}{2}
\setlist[parts,1]{label=(\alph*),leftmargin=*,itemsep=5pt,topsep=3pt}
\setlist[parts,2]{label=(\roman*),leftmargin=*,itemsep=3pt,topsep=2pt}

\newcommand{\pts}[1]{\hfill   \textit{[#1]}}


\begin{document}


\begin{center}
  {\large\bfseries BANARAS HINDU UNIVERSITY}\\[2pt]
  {\normalsize B.Com. (Hons.) F.M.M. --- Semester IV Examination, 2015-16}\\[6pt]
  \rule{0.8\linewidth}{0.5pt}\\[6pt]
  {\large\bfseries Paper No. BCF-299: Financial Analysis}\\[4pt]
  {\normalsize Time: 3 Hours\hfill Maximum Marks: 70}
\end{center}

\rule{\linewidth}{0.4pt}

\smallskip



\noindent   \textit{Instructions: Answer all questions. All questions carry equal marks (14 marks each).}

\bigskip




\noindent   \textbf{Question 1.} \\
What do you understand by financial statements analysis? Discuss the importance of such Statement.\pts{14}

\centerline{   \textbf{OR}}

The following information is obtained from Ganesh Co. Ltd. in a current year:\pts{14}

\begin{center}
\renewcommand{\arraystretch}{1.1}

\begin{tabular}{lr}
  Sales & Rs. 10,00,000 \\
  Variable Cost & Rs. 6,00,000 \\
  Fixed Cost & Rs. 3,00,000 \\
\end{tabular}
\end{center}
Find out the P/V ratio, BEP (Break Even Point) and MOS (Margin of Safety).

\medskip\rule{\linewidth}{0.2pt}




\noindent   \textbf{Question 2.} \\
Explain the meaning of ratio analysis. Discuss its importance in financial analysis.\pts{14}

\centerline{   \textbf{OR}}

Calculate the amount of Opening debtors and closing debtors from the following figures:\pts{14}

\begin{itemize}[itemsep=2pt,topsep=2pt]
  \item Debtors turnover ratio: 4 Times
  \item Cost of goods sold: Rs. 6,40,000
  \item Gross Profit ratio: 20\%
  \item Closing debtors were Rs. 20,000 more than at the beginning.
  \item Cash sales being $33\frac{1}{3}\%$ of credit sales.
\end{itemize}

\medskip\rule{\linewidth}{0.2pt}

\pagebreak




\noindent   \textbf{Question 3.} \\
What is funds flow statement? Give specimen of schedule of changes in the working capital and funds flow statement.\pts{14}

\centerline{   \textbf{OR}}

From the following figures prepare:\pts{14}

\begin{parts}
  \item A schedule of changes in working capital.
  \item Funds Flow statement.
\end{parts}


\begin{center}
\small
   \textbf{Balance Sheets} \\[4pt]
\renewcommand{\arraystretch}{1.2}

\begin{tabular}{|l|c|c|l|c|c|}
\hline
   \textbf{Capital \& Liabilities} &   \textbf{2014 (Rs.)} &   \textbf{2015 (Rs.)} &   \textbf{Properties \& Assets} &   \textbf{2014 (Rs.)} &   \textbf{2015 (Rs.)} \\ \hline
Equity share capital & 3,00,000 & 3,50,000 & Fixed Assets (Net) & 5,10,000 & 6,20,000 \\
Preference share capital & 2,00,000 & 1,00,000 & Investments & 30,000 & 80,000 \\
Profit \& Loss A/c & 1,10,000 & 2,70,000 & Current Assets & 2,40,000 & 3,75,000 \\
Debentures & 1,00,000 & 2,00,000 & Discount on Debentures & 10,000 & 5,000 \\
Current Liabilities & 80,000 & 1,60,000 & & & \\ \hline
   \textbf{Total} &   \textbf{7,90,000} &   \textbf{10,80,000} &   \textbf{Total} &   \textbf{7,90,000} &   \textbf{10,80,000} \\ \hline
\end{tabular}
\end{center}




\noindent   \textbf{Additional Information:}

\begin{enumerate}[label=(\roman*),itemsep=2pt,topsep=2pt]
  \item A Machine with book value of Rs. 40,000 was sold for Rs. 25,000.
  \item Depreciation charged during the year Rs. 60,000.
  \item Preference shares were redeemed at a premium of 15\%.
\end{enumerate}

\medskip\rule{\linewidth}{0.2pt}




\noindent   \textbf{Question 4.} \\
What do you mean by ``Cash Flow Statement''? How is it different from `Fund Flow Statement'? What is its utility?\pts{14}

\centerline{   \textbf{OR}}

From the following Balance Sheets, prepare Cash Flow Statement for the year ended 31st March, 2016:\pts{14}


\begin{center}
\small
   \textbf{Balance Sheets} \\[4pt]
\renewcommand{\arraystretch}{1.2}

\begin{tabular}{|l|r|r|l|r|r|}
\hline
   \textbf{Capital \& Liabilities} &   \textbf{1-04-2015} &   \textbf{31-03-2016} &   \textbf{Properties \& Assets} &   \textbf{1-04-2015} &   \textbf{31-03-2016} \\ \hline
Capital & 1,25,000 & 1,53,000 & Building & 35,000 & 60,000 \\
Loan from X & 25,000 & -- & Land & 35,000 & 50,000 \\
Loan from Bank & 40,000 & 50,000 & Machinery & 80,000 & 55,000 \\
Creditors & 40,000 & 44,000 & Stock & 40,000 & 25,000 \\
& & & Debtors & 30,000 & 50,000 \\
& & & Cash & 10,000 & 7,000 \\ \hline
   \textbf{Total} &   \textbf{2,30,000} &   \textbf{2,47,000} &   \textbf{Total} &   \textbf{2,30,000} &   \textbf{2,47,000} \\ \hline
\end{tabular}
\end{center}




\noindent   \textbf{Additional Information:} \\
During the year a machine costing Rs. 10,000 with accumulated depreciation Rs. 3,000 was sold for Rs. 5,000.

\medskip\rule{\linewidth}{0.2pt}

\pagebreak




\noindent   \textbf{Question 5.} \\
What is Leverage? Explain the operating leverage and financial leverage. How these are calculated?\pts{14}

\centerline{   \textbf{OR}}

A firm has sales of Rs. 25,00,000, variable cost of Rs. 17,50,000 and fixed cost Rs. 5,00,000 and debt of Rs. 2,50,000 at 10\% rate of interest. Compute operating, financial and combined leverages.\pts{14}

\vfill
\rule{\linewidth}{0.4pt}

\begin{center}\small
\href{https://bhu-exam.vercel.app/}{Click here to visit our website}\\[4pt]
\href{https://me-aryan.vercel.app/}{Click here to visit our contributor}\\[4pt]
\href{https://quashan-me.vercel.app/}{Click here to visit our contributor}
\end{center}

\end{document}
